\scanpage{75}\mspage{62}
% archive sheet number: 74

Si $\unsure{i'|a'|\alpha'}=i'_0$, $\alpha=\mathrm{id}_{i'_0}$, alors cette commutativité
revient à $\gamma=\beta$, \unsure{ça|ce} qui signifie qu'il
y a une flèche et une seule de
$\alpha$ dans $(i,\beta)$. Cela montre donc que
$(i'_0,\mathrm{id}_{i'_0})$ est un objet initial de $\mathord{}_{i'_0}\backslash C(u)$.
On a donc prouvé le lemme :

\uline{Lemme} 3.8. Soit \unsure{--} $u:I\unsure{\mapsto|\to} I'$ \deleted{\unsure{deux catégories}}\inserted{une flèche}
dans $\mathrm{Cat}$, factorisons-le en
\[
I\xrightarrow{w} C(u)\xrightarrow{v}I'\,,
\]
où $C(u)$ est le mapping-cone défini
p.ge 60, $w$ une immersion fermée
de complémentaire $I'$, $v$ la rétraction
canonique de $C(u)$ sur $I'$. Ceci
posé, pour \unsure{tt} $i'_0\in\mathrm{Ob}\, I'$, la catégorie
$\mathord{}_{i'_0}\backslash(C(\unsure{u|ui}),v)$ admet un objet initial ---
donc $v:C(u)\longrightarrow I'$ est (\unsure{tout ce|tant est} qu'il
y a de) \unsure{cofiltrante|cofiltrant}.

Nous pouvons maintenant prouver
3.7. Tout d'abord, les hypothèses
faites sur $I,I'$ se transportent à
$C(u)$ (hyp. artinienne, et \unsure{éventuelle|éventuelles}
hyp. \unsure{noethérienne|noethériennes} sur $\mathrm{Esp}(C(u))$. Donc
